\documentclass{article}
\usepackage[margin=1in]{geometry}
\usepackage{graphicx}
\usepackage{booktabs}
\usepackage{amsmath,amssymb}
\usepackage{hyperref}
\usepackage[backend=biber]{biblatex}
\usepackage{listings}

\title{Assignment 3 Report}
\author{
Isaac Levy \\
gbz6qn@virginia.edu}
\date{DS7400 - Deep Learning, Fall 2026}
\begin{document}

\section{Introduction}
\subsection{Background}
The goal of this project is to perform fingerprinting of stereoelectroencephalography (SEEG) electrode contacts. 
There are 3 patients, each of whom have 2 interictal (betwen seizure) recordings of ~5 minutes. 
I am training and testing a model for each individual patient, since the number of contacts and implantation
locations vary substantially between patients (these patients have 36, 49, and 77 good electrode contacts).

At this time, I am restricting my analyses to interictal recordings for three reasons:
\begin{enumerate}
    \item Interictal recordings are generally less noisy than ictal (seizure) recordings
    \item Interictal recordings are longer than ictal recordings, providing more training and testing data 
    \item I am interested in further exploring the clinical utility of interictal recordings for understanding seizure dynamics
\end{enumerate}

\subsection{Data Processing}
Preprocessing included various filtering and resampling steps.
First a 1 Hz high-pass filter was applied to remove slow drift and the DC offset that each electrode has,
which is intrinsic to the elctrode itself, rather than a biological signal, and is therefore not a feature we want to learn.
Next, a notch filter was applied at 60 Hz (and harmonic frequencies up to the Nyquist frequency) to remove power line noise.
Then, recordings were downsampled to 256 Hz from an original 500 Hz or 512 Hz to reduce computational overhead.
Finally, global average referencing was applied to all contacts (at the subect level) to remove artifacts.


Since each subject has two recordings of ~5 minutes, the first recording was split into training (180 seconds), 
validation (55 seconds), and test (55 seconds) sets with 5 second gaps in between the sets, while 
the second recording was only used for testing. To create more training samples, the 1-second window was 
moved 0.1 seconds (26 data points at 256 Hz) forwards to generate overlapping windows, 
increasing the training size from 180 to 1,700+ sammples.
Additionally, global normalization was applied:
the mean and standard deviation was computed over every 1-second window and conctact in the training set, and
the training data, validation data, test data, and run2 test data were z-scored using these values.
This rescales every window while preserving amplitude differences between contacts and only uses information computed
from the training set.

\section{Baseline Models}
\subsection{Data Processing}
I created two baseline models for this fingerprinting task:
\begin{enumerate}
    \item A fully-connected MLP with 2 hidden layers
    \item A 1D-CNN with a transformer encoder (code sourced from CITE)
\end{enumerate}

PUT IN DIAGRAM OF CNN

Both models were designed to directly use the 1-second long snippets of the interictal recordings.
This data representation was originally inspired by CITE LILY, which used accelerometry data from a wrist sensor and a lag-based
feature extractor to fingeprint individuals.

\subsection{Results}


\section{Proposed Model}
\subsection{Data Representation}
\subsection{Model Architecture}
\subsection{Hypotheses}

\section{Results}



\end{document}